% Cuestión VII. Toda cifra medida sale de generado/ (\dato o tablas \input); no escribir cifras a mano.

\subsection{El mecanismo}

LSB común reemplaza el bit menos significativo de cada byte que usa por un bit del mensaje. Si los bits del
mensaje y los bits bajos del portador no guardan relación, cada byte usado cambia con probabilidad un medio, de modo que
se modifica alrededor de la mitad de los bytes que llevan mensaje. Majeed y Sulaiman~\cite{majeed2015} aprovechan una propiedad
que el propio documento no enuncia: LSB sólo toca el bit 0 de cada byte, así que los bits 1 y 2 son idénticos en el
portador y en el archivo estego. Con esos dos bits se arman cuatro grupos (patrones \texttt{00}, \texttt{01}, \texttt{10} y \texttt{11}) que el emisor
y el receptor calculan de la misma manera, sin necesidad de intercambiar nada.

Para cada grupo, el emisor cuenta cuántos de sus bytes con mensaje cambiarían con LSB común (sea $c$) y cuántos quedarían
igual ($n-c$, con $n$ el total del grupo). Si cambiarían más de los que no ($c > n-c$), en ese grupo escribe el
complemento de los bits del mensaje, de modo que ahora cambian justo los $n-c$ bytes que antes no cambiaban, y deja la bandera
del grupo en uno. El receptor lee las cuatro banderas, recalcula el patrón de cada byte y aplica la bandera de su grupo para
recuperar el bit original. Por grupo, la cantidad final de bytes modificados es
\[
  \min(c,\; n-c) \le \min\!\left(c,\; \tfrac{n}{2}\right),
\]
donde un empate no se invierte y deja $c = n-c$. Sumando los cuatro grupos, LSBI nunca modifica más bytes con mensaje que LSB común sobre las
mismas posiciones, y a lo sumo modifica la mitad de los bytes de cada grupo. El costo es de cuatro bytes de banderas, escritos con
LSB1, que pueden agregar como mucho cuatro cambios. La ganancia en cada grupo es $|c - n/2|$, es decir, cuánto se aparta el conteo del
punto medio: cuanto más desbalanceado sea el grupo, más se gana.

\subsection{Evidencia medida}

El cuadro~\ref{tab:q7-inversion} compara, sobre los mismos bytes, los cambios sin inversión (lo que haría LSB común en esas posiciones)
con los cambios con inversión, para cinco cargas: el PNG de la cátedra, el mensaje de prueba del paper, datos aleatorios y los
dos vectores cifrados de la cátedra. La última columna cuenta los cambios que producen las banderas mismas.

\begin{table}[htbp]
\centering
\footnotesize
\setlength{\tabcolsep}{3pt}
\resizebox{\linewidth}{!}{\begin{tabular}{lrrrrcr}
\toprule
Caso & Bytes con mensaje & Sin inversión & Con inversión & Reducción (\%) & Banderas & Cambios en banderas \\
\midrule
PNG (sin encripción) & 359\,160 & 188\,308 & 170\,852 & 9{,}27 & \texttt{1111} & 0 \\
Mensaje del paper & 1\,880 & 1\,044 & 832 & 20{,}31 & \texttt{0001} & 3 \\
Datos aleatorios & 320\,072 & 160\,032 & 159\,771 & 0{,}16 & \texttt{0101} & 2 \\
PNG aes256-ofb (cátedra) & 359\,192 & 179\,491 & 179\,329 & 0{,}09 & \texttt{0001} & 3 \\
PNG 3des-cfb (cátedra) & 359\,192 & 179\,425 & 179\,195 & 0{,}13 & \texttt{1000} & 3 \\
\bottomrule
\end{tabular}
}
\caption{Bytes con mensaje modificados sin y con inversión de patrones, para cinco cargas ocultadas en \texttt{lado.bmp} con LSBI. La reducción se calcula sobre los bytes con mensaje; las banderas se cuentan aparte}
\label{tab:q7-inversion}
\end{table}

Con el PNG, la inversión baja los cambios de \dato{inv-png-sin} a \dato{inv-png-con}, una reducción del \dato{inv-png-reduccion}\,\%, y las
banderas no agregan ningún cambio. El cuadro~\ref{tab:q7-patrones} muestra de dónde sale esa ganancia.

\begin{table}[htbp]
\centering
\small
\begin{tabular}{crrcr}
\toprule
Patrón & Cambiados & Sin cambio & Invertido & Cambios finales \\
\midrule
\texttt{00} & 32\,367 & 32\,028 & sí & 32\,028 \\
\texttt{01} & 32\,650 & 32\,263 & sí & 32\,263 \\
\texttt{10} & 35\,119 & 34\,980 & sí & 34\,980 \\
\texttt{11} & 88\,172 & 71\,581 & sí & 71\,581 \\
\midrule
Total & 188\,308 & 170\,852 & & 170\,852 \\
\bottomrule
\end{tabular}

\caption{Conteo por patrón al ocultar el PNG de la cátedra con LSBI: bytes que cambiarían, bytes que quedarían igual, si el grupo se invierte y cambios que quedan después}
\label{tab:q7-patrones}
\end{table}

Los cuatro grupos se invierten, pero en tres de ellos la diferencia entre cambiados y sin cambio es mínima, y casi toda la reducción sale del
patrón \texttt{11}, que es el más numeroso y el de mayor desequilibrio. La explicación es una propiedad del portador y no del método: en
\texttt{lado.bmp} hay muchas zonas saturadas, cuyos bytes valen 255, tienen el patrón \texttt{11} y el bit menos significativo en uno; en
ese grupo los bits bajos del portador son, por lo tanto, mayoritariamente unos (un conteo auxiliar sobre \texttt{lado.bmp} lo confirma), mientras que el PNG,
un archivo comprimido, tiene menos unos que ceros; así, en ese grupo cambian más bytes de los que quedan. Es coherente con la proporción global de unos del canal
azul del portador (\dato{est-portador-B-lsb}, mayor que un medio). En otras palabras, LSBI gana cuando el portador y el mensaje tienen estadísticas distintas dentro de algún
patrón, y esa diferencia es propia del par portador-mensaje.

Los casos cifrados y el aleatorio confirman esto. Con datos aleatorios y con los dos vectores cifrados de la cátedra (AES-256 en OFB y 3DES en
CFB) la reducción es de \dato{inv-azar-reduccion}\,\%, \dato{inv-aes256ofb-reduccion}\,\% y \dato{inv-descfb-reduccion}\,\%: mínima. Un mensaje cifrado se comporta
como una secuencia de bits equiprobables e independientes del portador. En cada grupo el conteo $c$ es entonces una variable
binomial de media $n/2$, y la ganancia esperada es su desviación absoluta media, aproximadamente $\sqrt{n/(2\pi)}$ por grupo. Para la carga aleatoria eso
suma unos \dato{inv-azar-teorico} cambios, el \dato{inv-azar-teoricopct}\,\% de los bytes con mensaje. Lo medido queda por debajo pero
del mismo orden, con la dispersión propia de una única muestra. La ganancia proviene del azar del muestreo y no de una regularidad que el método
aproveche; además, con las banderas, la ganancia neta baja todavía más (hasta \dato{inv-aes256ofb-banderas} cambios en banderas en los vectores cifrados).

La comparación con LSB1 sobre el mismo PNG (\hyperref[sec:cuestion-ii]{Cuestión II}) va en el mismo sentido. LSB1 usó \dato{png-LSB1-usados} bytes y cambió
\dato{png-LSB1-cambiados}, una tasa del \dato{png-LSB1-tasa}\,\%; LSBI usó \dato{png-LSBI-usados} bytes (contando las banderas) y cambió \dato{png-LSBI-cambiados}, una tasa
del \dato{png-LSBI-tasa}\,\%. El PSNR pasa de \dato{png-LSB1-psnr} a \dato{png-LSBI-psnr}\,dB: una mejora de una fracción de decibel, coherente con una
reducción de cambios de menos del diez por ciento. Con el mensaje corto del paper los cambios pasan de \dato{paper-LSB1-cambiados} a \dato{paper-LSBI-cambiados}, y el PSNR de
\dato{paper-LSB1-psnr} a \dato{paper-LSBI-psnr}\,dB, pero en ese régimen ambas cifras son muy altas y no dicen mucho sobre la calidad del método. La reducción relativa
del texto del paper es la mayor de las cinco cargas (\dato{inv-paper-reduccion}\,\%), y es esperable: un texto ASCII tiene siempre en cero el bit más significativo de cada carácter y sus
bits están lejos de ser equiprobables. Pero se calcula sobre \dato{inv-paper-n} bytes solamente y con \dato{inv-paper-banderas} cambios de banderas que hay que sumar aparte, así que la
ganancia neta es algo menor.

\subsection{¿Es realmente una mejora?}

En sentido estricto y medible, sí: LSBI no modifica más bytes con mensaje que LSB común sobre las mismas posiciones (a lo sumo cuatro cambios más por las banderas),
y cuando el mensaje tiene una estadística distinta de la del portador modifica menos, con un PSNR algo mayor. El precio y los límites de esa mejora son estos.

\textbf{La capacidad baja.} Como los bits de mensaje sólo van a los bytes verde y azul, y además hay que reservar las banderas, el mismo
\texttt{lado.bmp} admite \dato{cap-LSBI} bytes con LSBI contra \dato{cap-LSB1} con LSB1. El documento dice tener en cuenta la capacidad (p.~343), pero no la mide.

\textbf{La ganancia depende del mensaje.} Justamente en el caso que más importa para este programa, un mensaje cifrado, la reducción es casi nula
(cuadro~\ref{tab:q7-inversion}), y el algoritmo es más complejo que LSB común: dos pasadas y un conteo por patrón.

\textbf{No vuelve indetectable el mensaje.} El estadístico $\chi^2$ de pares de valores de Westfeld y Pfitzmann~\cite{westfeld2000}, medido en el
cuadro~\ref{tab:q2-estadisticas}, baja en el canal azul de \dato{est-portador-B-chi} a \dato{est-png-LSBI-B-chi} y en el verde de \dato{est-portador-G-chi} a
\dato{est-png-LSBI-G-chi}: LSBI sigue siendo un reemplazo de bits menos significativos, y un mensaje casi aleatorio iguala los conteos de cada par de valores
haya inversión o no. El análisis RS de Fridrich, Goljan y Du~\cite{fridrich2001} apunta a la misma asimetría del reemplazo de LSB; no lo aplicamos, y no pretendemos
que LSBI lo evada ni que lo detecte peor. El canal rojo, además, queda con el mismo estadístico del portador (\dato{est-png-LSBI-R-chi}), así que la diferencia entre
canales es una marca por sí misma, visible en la figura~\ref{fig:q2-mapas}.

\textbf{La seguridad adicional es oscuridad.} Que el rojo ``actúe como ruido'' sólo protege mientras el método sea secreto; quien lo conoce sabe qué canales
mirar, contra el principio de Kerckhoffs~\cite{kerckhoffs1883}. El mensaje, además, empieza siempre en el mismo byte y se recorre en orden secuencial, con lo que la extracción es
trivial para cualquiera que conozca el método. La inversión de patrones no agrega ninguna clave.

En conclusión, consideramos que LSBI es una mejora real pero estrecha. Reduce la cantidad de bytes modificados, y eleva algo el PSNR, cuando la estadística
del mensaje difiere de la del portador dentro de algún patrón. A cambio pierde un tercio de la capacidad más cuatro bytes, y su ventaja se diluye cuando el mensaje se cifra, que es la práctica recomendable. No mejora la resistencia
al estegoanálisis estadístico ni justifica, con lo medido, decir que el mensaje pasa inadvertido: los rastros estadísticos de los canales verde y azul, y la asimetría con el rojo,
siguen presentes.
